\section{Norm balls and weighted power cones}
\label{sec:lp-power}
The local bound depends on a curved boundary patch, rather than on
rotational symmetry. It therefore gives the same unrestricted dimension
frontier for every $\ell_p$ ball with $1<p<\infty$. The global frontier
also persists, but its proof uses the primal and polar boundaries as
independent $C^1$ manifolds. This distinction matters at coordinate zeros.

\subsection{Exact counts and the dimensions of count-optimal lifts}
Write
\[
 B_p^N=\{x\in\R^N:\textstyle\sum_j|x_j|^p\leq1\},\qquad
 K_{p,m}=\{(t,u)\in\R\times\R^{m-1}:t\geq\norm{u}_p\}.
\]
Fix an o-minimal expansion in which the relevant fixed powers are
definable; $\R_{\exp}$ suffices. For rational exponents one can work
entirely with semialgebraic cones and lifts.

\begin{theorem}[Dimension-capped $\ell_p$ balls]\label{thm:lp-granularity}
Let $1<p<\infty$, $N\geq2$, and let $d\geq3$ be an integer.
Among exact lifts of $B_p^N$
through arbitrary definable proper cones of dimensions at most $d$, with
arbitrary affine equations, projections, and free variables, the minimum
number of factors and minimum total cone dimension are simultaneously
\begin{equation}\label{eq:lp-generic-frontier}
 k_* = \left\lceil\frac{N-1}{d-2}\right\rceil,
 \qquad M_* = N-1+2k_*.
\end{equation}
Every such lift has at least $k_*$ factors of dimension at least three,
even if arbitrarily many rays are allowed.
\end{theorem}
\begin{proof}
On the boundary patch where every coordinate is nonzero,
$f(x)=\sum_j|x_j|^p$ is smooth and
\[
 D^2 f(x)=p(p-1)\operatorname{diag}(|x_j|^{p-2})\succ0.
\]
Its restriction to the tangent space, divided by $\norm{Df(x)}$, is
the positive second fundamental form. After the reductions of
Lemma~\ref{lem:reduction}, Theorem~\ref{thm:curvature} gives
$N-1\leq\sum_i(s_i-2)_+$, where $s_i\leq d$ are the reduced factor
dimensions. If $k_+$ factors have $s_i\geq3$, then
$k_+\geq k_*$ and their dimensions sum to at least $N-1+2k_+$.
Original dimensions and the total factor count are no smaller.

For attainment partition $N-1=b_1+\cdots+b_{k_*}$ with
$1\leq b_i\leq d-2$. Make a chain of internal tree nodes: node $i<k_*$
has $b_i$ new leaf children and node $i+1$; the last node has $b_{k_*}+1$
leaf children. Node $i$ requires its scalar to dominate the $p$-norm of
its children, using $K_{p,b_i+2}$, and the root scalar is one. There are
$N$ leaves. Induction from the leaves shows that a node scalar dominates
the $p$-norm of its descendant leaves. Conversely, assigning exactly
those subtree norms makes every internal constraint feasible whenever
$\norm{x}_p\leq1$. The projection is therefore exactly $B_p^N$, and
the total dimension is $\sum_i(b_i+2)=N-1+2k_*$. At zero leaves choose
strictly positive internal scalars decreasing from the root to obtain
Slater feasibility.

The cones $K_{p,m}$ are closed, pointed, full-dimensional and convex.
For $p=a/b$ rational, nonnegative auxiliary variables satisfying
$z_j^b=|u_j|^a$, $w^b=t^a$, $t\geq0$ and $\sum_j z_j\leq w$
give a semialgebraic description after projection. This also handles an
even denominator without changing the meaning of the power.
\end{proof}

The tree construction is established prior work: the three-dimensional
version is \citet[Proposition~1]{KrokhmalSoberanis2010}; see also
\citet{BlancoMartinezAnton2024} for representations by specified cone
families. The contribution of Theorem~\ref{thm:lp-granularity} is its
matching lower bound against arbitrary definable proper factors and
affine lifts, together with the total-dimension optimum. The polyhedral
endpoints $p=1,\infty$ have no such curved patch and are excluded.
Restriction to interior two-dimensional sections of the positive-capacity
factors also gives the ambient barrier lower bound $2k_*$, as in
Corollary~\ref{cor:proper-smooth-frontier}. Its attainment is asserted
here only for $p=2$.

\begin{proposition}[All capacity excesses at the minimum count]
\label{prop:lp-capacity-excess}
Under the hypotheses of Theorem~\ref{thm:lp-granularity}, set $n=N-1$,
$c=d-2$, and $\sigma=k_*c-n$. Suppose an exact lift uses
exactly $k_*$ original factors, of dimensions $m_i\leq d$, and let $s_i$
be their dimensions after minimal-face reduction. Then every $s_i\geq3$
and
\begin{equation}\label{eq:lp-excess}
 \Delta=\sum_i(s_i-2)-n\in\{0,\ldots,\sigma\},\qquad
 \sum_i s_i=n+2k_*+\Delta.
\end{equation}
Every indicated $\Delta$ occurs in a full-face $p$-norm tree lift.
If $c$ divides $n$, every count-optimal lift has $s_i=m_i=d$ for all $i$.
\end{proposition}
\begin{proof}
If one reduced factor had dimension at most two, the remaining capacity
would be at most $(k_*-1)c<n$, a contradiction. The lower curvature
bound and the dimension cap now give \eqref{eq:lp-excess}. For any desired
$\Delta$, start with the preceding chain and choose integers
$0\leq r_i\leq c-b_i$ with $\sum_i r_i=\Delta$. Replace
$K_{p,b_i+2}$ by $K_{p,b_i+2+r_i}$ and fix the added norm coordinates
to zero. This preserves the projection and strict feasibility, so the
larger displayed cones remain their full minimal faces. If $\sigma=0$,
equality in $\sum_i(s_i-2)\leq k_*c$ forces every $s_i=m_i=d$.
\end{proof}

At $\Delta=0$ the local conclusion is stronger than the arithmetic.
On a patch where every coordinate is nonzero the supporting-polarity map
is smooth, so the generic $C^2$ selections can be written on a common
contact chart. Put $H_i=-DA_i^*DB_i$ for those selections
$A_i(u),B_i(v)$. Each kernel
$\ip{A_i(u)}{B_i(v)}$ is nonnegative and zero when $u=v$. Its Hessian
at a diagonal point is positive semidefinite and annihilates $(h,h)$;
thus $H_i$ is symmetric positive semidefinite. The total
$H=\sum_iH_i$ is positive definite. Rank equality gives
$\rank H_i=s_i-2$, and the linear-algebra argument in
Section~\ref{sec:regularity} gives
\begin{equation}\label{eq:lp-local-projections}
 P_i=H^{-1/2}H_iH^{-1/2},\qquad
 P_i^2=P_i,\qquad P_iP_j=0\ (i\ne j),\qquad \sum_iP_i=I.
\end{equation}
Moreover, both quotient derivatives into
$B_i^\perp/\R A_i$ and $A_i^\perp/\R B_i$ are onto: otherwise the
mixed rank would be smaller than $s_i-2$. These are local statements;
they do not determine the full cones.

\begin{example}[A saturated lift with a non-strictly-convex factor]
\label{ex:lp-halfcone}
For $B_p^3$, set
\[
 K_1=K_{p,3},\qquad
 K'_2=\{(s,t,z)\in K_{p,3}:t\geq0\}.
\]
The two constraints
\[
 (t,x_1,x_2)\in K_1,\qquad (1,t,x_3)\in K'_2
\]
project exactly onto $B_p^3$: the first already implies $t\geq0$, so
the extra halfspace in the second is redundant on the affine lift.
At $x=0$ and $0<t<1$ both constraints are strict. Thus both full
three-dimensional factors are necessary in the curvature count, and
their total capacity is $2=3-1$. Nevertheless $K'_2$ has the
two-dimensional exposed face
$\{(s,0,z):s\geq|z|\}$, whereas every proper nonzero face of $K_{p,3}$
is a ray. The factors are not linearly isomorphic. At every contact
with $(x_1,x_2)\ne0$, the added facet is inactive and the local cone
germ is unchanged. In particular, saturation alone cannot classify
the global factors or justify a target-dependent barrier premium.
\end{example}

\subsection{Global selections and the Young construction}
Let $q=p/(p-1)$. The two boundaries $\partial B_p^N$ and
$\partial B_q^N$ are $C^1$ and strictly convex. A global factorization
in this subsection means $C^1$ primal selections on the first boundary
and genuine $C^1$ normalized dual selections on the second, satisfying
the full two-point identity for $1-\ip xy$. The two selections are
differentiated in their own boundary coordinates.

\begin{theorem}[Global $C^1$ frontier for $\ell_p$ balls]
\label{thm:lp-smooth-frontier}
Let $N\geq3$ and let $d\geq3$ be an integer. Among exact lifts with
these global selections through arbitrary proper factors of dimension
at most $d$, the simultaneous
minimum capacity $S=\sum_i(\dim K_i-2)_+$, total factor count $k$, and
total cone dimension $M$ are
\begin{equation}\label{eq:lp-smooth-frontier}
 (S,k,M)=
 \begin{cases}
 (N,h,N+2h),\quad h=\lceil N/(d-2)\rceil,&d<N+1,\\
 (N-1,1,N+1),&d\geq N+1.
 \end{cases}
\end{equation}
The lower bound does not require definability of the cones or selections.
\end{theorem}
\begin{proof}
At a contact $\ip xy=1$ the tangent spaces are $y^\perp$ and
$x^\perp$. Their Euclidean pairing is nondegenerate, since its first
kernel is $y^\perp\cap\R x=\{0\}$. Consequently $S\geq N-1$ by
Lemma~\ref{lem:universal-channel}. If equality holds,
Theorem~\ref{thm:general-joint-cover} allows exactly one positive-capacity
factor, of dimension $N+1$: both boundaries are $C^1$ diffeomorphic to
$\Sph^{N-1}$. A strict cap excludes this case, so integrality gives
$S\geq N$. Counting positive-capacity factors then gives $k\geq h$
and $M\geq N+2h$. For $d\geq N+1$, the direct lift $K_{p,N+1}$ with
primal and dual maps $(1,x),(1,-y)$ attains the second line.

For the first line define, for $r\geq1$,
\begin{equation}\label{eq:grouped-p-cone}
 \mathcal P_{p,r}
 =\{(u,v,w)\in\R_+^2\times\R^r:
                uv^{p-1}\geq\norm{w}_p^p\}.
\end{equation}
This is the closed perspective epigraph of $\norm{w}_p^p$, hence a
proper cone of dimension $r+2$. At $v=0$ the definition gives $w=0$
and $u\geq0$, exactly the closure of its part with $v>0$.
Partition the coordinates into $h$ nonempty groups $G$ of sizes at most
$d-2$, and impose
\begin{equation}\label{eq:grouped-p-lift}
 (u_G,v_G,w_G)\in\mathcal P_{p,|G|},\qquad
 v_G=1,\qquad \sum_Gu_G=1,\qquad x_G=w_G.
\end{equation}
The projection is $B_p^N$, and $x=0$, $u_G=1/h$ is a Slater point.
On the boundary the primal fiber is unique, with $u_G=\norm{x_G}_p^p$.
Define
\begin{equation}\label{eq:young-group-factors}
 A_G(x)=(\norm{x_G}_p^p,1,x_G),\qquad
 B_G(y)=\left(\frac1p,\frac{\norm{y_G}_q^q}{q},-y_G\right).
\end{equation}
For $v>0$, Young's inequality gives
\[
 \frac{u}{p}+\frac{v\norm{y_G}_q^q}{q}-\ip{w}{y_G}
 \geq v\left(\frac{\norm{w/v}_p^p}{p}
             +\frac{\norm{y_G}_q^q}{q}-\ip{w/v}{y_G}\right)\geq0.
\]
Continuity includes $v=0$, so $B_G(y)$ belongs to the dual cone.
Its coefficients are the multipliers $1/p$ for $\sum_Gu_G=1$ and
$\norm{y_G}_q^q/q$ for $v_G=1$, minus the objective on $w_G$.
Their objective value is $1/p+\sum_G\norm{y_G}_q^q/q=1$.
Thus these are genuine certificates on the whole affine slice.
They give
\[
 \sum_G\ip{A_G(x)}{B_G(y)}=1-\ip xy
 \quad(x\in\partial B_p^N,\ y\in\partial B_q^N).
\]
Both maps are globally $C^1$ because $p,q>1$. Their total capacity,
count and dimension are $N,h,N+2h$, respectively.
\end{proof}

The scalar case $r=1$ of \eqref{eq:grouped-p-cone} is the power cone
$P_{1/p}=\{(u,v,w):u,v\geq0,\ u^{1/p}v^{1/q}\geq|w|\}$.
The corresponding standard $p$-norm model is recorded in
\citet[Section~4.2.2]{MosekPowerCookbook}. The construction and Young's
inequality are not novelty claims. Theorem~\ref{thm:lp-smooth-frontier}
proves their optimal resource count under the stated global regularity
condition. For $d=3$ it gives exactly $N$ factors, including when the
dictionary is restricted to three-dimensional power cones.

The scalar selections in \eqref{eq:young-group-factors} are both $C^2$
only when $p=2$. If $1<p<2$, the primal map fails to be $C^2$ at
coordinate zeros and the dual map is $C^2$; for $p>2$ the roles reverse.
This does not affect the preceding proof, which does not differentiate
the polarity map $y_j=\operatorname{sign}(x_j)|x_j|^{p-1}$.
Every ambient barrier has parameter at least $2h$ in the strict-cap
case, by the product orthant-section argument of
Corollary~\ref{cor:proper-smooth-frontier}. This is attained for $p=2$;
no such equality is asserted for the nonsymmetric perspective cones.
The reduced affine domain in \eqref{eq:grouped-p-lift} also has the
intrinsic barrier lower bound $h$: fix all allocations $u_G$ positive
with sum one, set all but one coordinate per group to zero, and vary
the remaining coordinate in
$(-u_G^{1/p},u_G^{1/p})$. This gives an $h$-dimensional open box through
the domain interior, whose boundary lies on the domain boundary.
Restriction and the box lower bound prove the assertion. The Euclidean
case is attained by the grouped barrier of Section~\ref{sec:concrete-barriers};
the construction here makes no matching intrinsic-parameter assertion
for $p\ne2$.

\begin{proposition}[The two-dimensional power-cone exception]
\label{prop:lp-two-dimensional-power}
For $N=2$, one arbitrary three-dimensional proper cone suffices for
$B_p^2$, with global $C^1$ primal and dual selections. Within the
dictionary of three-dimensional power cones, the minimum total factor
count is one for $p=2$ and two for $p\ne2$, even without requiring
global selections.
\end{proposition}
\begin{proof}
The direct $K_{p,3}$ lift proves the first assertion, and the two scalar
Young factors give the upper bound in the power dictionary. A one-factor
power-cone lift has total dimension three; by
Theorem~\ref{thm:minimal-dimension-rigidity} its cone would be linearly
isomorphic to $K_{p,3}$. We show that this is impossible unless $p=2$.

For $P_\alpha$, $0<\alpha<1$, the only projective boundary rays with
possibly non-smooth second derivatives are its two axial rays. In the
chart $u=1$ at the first, the boundary is
$v=|w|^{1/(1-\alpha)}$; in the chart $v=1$ at the second it is
$u=|w|^{1/\alpha}$. Elsewhere the boundary is smooth with positive
curvature. If $\alpha\ne1/2$, exactly one of these exponents is below
two, so exactly one boundary ray is not $C^2$. For $K_{p,3}$ the
projective section is the $p$-unit circle. Near $(1,0)$ it has graph
\[
 x_1=(1-|x_2|^p)^{1/p}
     =1-\frac{|x_2|^p}{p}+o(|x_2|^p),
\]
and the same holds at its other three coordinate points. It therefore
has exactly four non-$C^2$ rays when $p<2$, and none when $p>2$.
An invertible linear cone map induces smooth projective coordinate
changes and preserves this count. This excludes $\alpha\ne1/2$.
For $\alpha=1/2$, the section is a nonsingular conic with nonzero
second fundamental form everywhere. When $p>2$ the four coordinate
points instead have zero second fundamental form; its vanishing is
preserved by projective transformations, as follows by differentiating
a homogeneous supporting equation along its tangent. When $p<2$ the
$C^2$ obstruction already suffices. At $p=2$ both cones are Lorentz
cones, which proves the remaining case.
\end{proof}

\subsection{A curvature bound for generalized Euclidean power cones}
Let $\alpha_j>0$, $\sum_{j=1}^m\alpha_j=1$, and set
\[
 g(x)=\prod_{j=1}^m x_j^{\alpha_j},\qquad
 C_{\alpha,k}=\{(x,z)\in\R_+^m\times\R^k:
                       \norm{z}_2\leq g(x)\}.
\]
The continuous extension of $g$ to the nonnegative orthant is understood.
It is concave, so $C_{\alpha,k}$ is a proper convex cone. It is
definable in $\R_{\exp}$, and semialgebraic for rational weights.

\begin{theorem}[Power-cone curvature capacity]\label{thm:power-capacity}
Let $m,k\geq1$ and $m+k\geq3$. Every exact affine lift of
$C_{\alpha,k}$ by definable proper cones of dimensions $d_i$ satisfies
\begin{equation}\label{eq:power-capacity}
 \sum_i(d_i-2)_+\geq m+k-2.
\end{equation}
In particular, a lift using three-dimensional Lorentz cones and any
number of rays requires at least $m+k-2$ Lorentz factors. The same
scalar bound $m-1$ holds for the nonnegative geometric-mean cone
$\{(x,t):x\geq0, 0\leq t\leq g(x)\}$ when $m\geq2$.
\end{theorem}
\begin{proof}
Intersect the cone with $\sum_jx_j=1$. This is a compact convex body
of dimension $m+k-1$ in its affine hull: $x$ ranges over the simplex
and $\norm z\leq g(x)$ is uniformly bounded. Restricting an exact lift
by this affine equation introduces no cone factors. On the patch
$x_j>0$, $z\ne0$, the boundary is $\Phi(x,z)=0$, where
$\Phi(x,z)=\norm z-g(x)$. Write $z=g(x)\omega$ with
$\norm\omega=1$. A tangent vector $(h,\xi)$ obeys
\[
 \sum_jh_j=0,\qquad
 \ip\omega\xi=Dg(x)[h].
\]
The Hessian on such tangent vectors is
\begin{equation}\label{eq:power-second-form}
 D^2\Phi[(h,\xi),(h,\xi)]
 =-D^2g(x)[h,h]
  +\frac{\norm{\xi}_2^2-\ip\omega\xi^2}{g(x)}.
\end{equation}
Direct differentiation gives the weighted variance identity
\[
 -D^2g(x)[h,h]
 =g(x)\left\{\sum_j\alpha_j\left(\frac{h_j}{x_j}\right)^2
       -\left(\sum_j\alpha_j\frac{h_j}{x_j}\right)^2\right\}.
\]
It is nonnegative and vanishes exactly when $h$ is proportional to
$x$. On $\sum_jh_j=0$, this forces $h=0$. If also the second term
of \eqref{eq:power-second-form} vanishes, then $\xi$ is proportional
to $\omega$; the tangent equation with $h=0$ forces $\xi=0$.
Thus the boundary patch has strictly positive second fundamental form
of rank $m+k-2$. Theorem~\ref{thm:curvature}, applied in the affine
hull of this compact section, gives \eqref{eq:power-capacity}.
Minimal-face reduction can only reduce the displayed dimensions.
For the nonnegative scalar cone the positive upper graph $t=g(x)$
has exactly the same Hessian on the simplex tangent; its section has
dimension $m$ and curved boundary dimension $m-1$. The case $m=1$
in the vector statement reduces to the ordinary Euclidean ball section.
\end{proof}

This result is a dimension bound, not an exact count for every rational
weight. \citet[Theorem~3.6]{Wang2024} proves the scalar bound $m-1$ for
simple mediated-set representations, together with separate
denominator-dependent bounds; his Theorem~3.15 gives $\lceil m/2\rceil$
for arbitrary three-dimensional SOC lifts. The proof above permits
arbitrary affine cone-coordinate maps, auxiliary equations and free
variables, and also gives the Euclidean-output term $k-1$.
\citet[Definition~1 and Theorem~14]{BlancoMartinezAnton2024} relate
extended representations and mediated graphs; their constructive maps
select triples of mediated variables. We make no priority claim for
the scalar inequality from this comparison alone, and do not extend
denominator-sensitive bounds beyond their established representation
classes. Theorem~\ref{thm:power-capacity} supplies a self-contained
geometric proof in the affine-lift model used throughout this paper.
